\section{What the model still lacks, and what it would change}\label{sec:limits}
The following omissions are the ones that can change an interpretation in
this paper. For each, the measured size of the effect is given where one
exists.

\paragraph{Atmospheres below 4600 K, above $\log g=6.1$, and without hydrogen.}
Track S stops where its TLUSTY families stop, and the flash continuations
rely on a hydrogen table whose normalization differs from TLUSTY by 49\% in
joining pressure (Section~\ref{sec:flash}). Cooling through 1000, 500 and
100 K requires dense, cool atmospheres with nonideal chemistry,
pressure-dependent and collision-induced absorption, and refraction
\cite{blouin18}. A hydrogen-free atmosphere is not yet possible: the source
programs normalize abundances to hydrogen, and a test at 4017 K shows that
rescaling all number abundances by one factor preserves composition but can
change individual absorption coefficients by a factor of 27.3. These gaps
decide whether a computed cooling age exists at all; nothing in
Sections~\ref{sec:hydrogen}--\ref{sec:flash} is a cooling age.

\paragraph{Grains.}
Equilibrium chemistry finds no condensates in the 3000--3200 K atmospheres
tested nor in the hydrogen-poor structures through 6400 K, but finds them in
the cool upper layers of 2600--2800 K models; dusty models become relevant
near 2800 K and below \cite{chabrier00}. Removing the condensed species from
the gas at 2800 K changes the joining pressure by 0.1246\% and temperature by
$-0.02347\%$. A grain-opacity calculation on one fixed 2800 K profile, with
a condensed mass fraction of $1.094\times10^{-4}$ in $0.1\,\mu$m compact
grains, gives integrated optical depths near $1\,\mu$m of 0.002610
(absorption) and 0.002829 (scattering); coupling grains to the atmospheric
structure and thermodynamics has not been done. Grains are a possible
correction during contraction, the cool initial main sequence and the
eventual cold-remnant phase. The continuous model reaches the main sequence
near 2765 K, so the grain-free assumption also needs assessment there.

\paragraph{Carbon conversion and molecular opacity.}
A fixed metal mixture misrepresents the cool opacity once carbon has been
converted to nitrogen, because CO formation couples the available carbon and
oxygen \cite{lederer09}. Complete conversion of the GS98 carbon captures a
hydrogen mass fraction of $5.728\times10^{-4}$. At $X=0.20$ and
$\rho=5.995\times10^{-6}$ g cm$^{-3}$ the Rosseland absorption mean changes
by the ratios in Table~\ref{tab:cn}; across densities of
$10^{-8}$--$10^{-4}$ g cm$^{-3}$, half conversion raises the 2705 K mean by
18.67--30.22\%. In a coupled 3200 K atmosphere at $\log g=5.15$, half
conversion raises the joining temperature by 1.601\% and lowers the pressure
by 5.681\%; applied to the $3.400\Tyr$ structure for 10 Myr this changes
the luminosity by $-5.162\%$ and the effective temperature by $-1.361\%$. At
5400 K the response is negligible. The effect is material for the cool
part of the track and must be included in a track that carries converted
catalysts to the surface.

\begin{table}[H]\centering\small
\caption{Rosseland absorption mean relative to the fixed GS98 mixture at
$X=0.20$ and $\rho=5.995\times10^{-6}$ g cm$^{-3}$ (gas only).}\label{tab:cn}
\begin{tabular}{@{}rrr@{}}\toprule
$T$ (K) & Half the carbon converted & All the carbon converted \\\midrule
2705 & 1.225 & 1.407 \\
3521 & 1.151 & 1.263 \\
4583 & 0.9971 & 0.9912 \\
5964 & 0.9901 & 0.9744 \\\bottomrule
\end{tabular}
\end{table}

\paragraph{Microscopic transport in degenerate and partially ionized matter.}
At the centre of the $3.890\Tyr$ Track P model the temperature is 0.1506
of the electron Fermi temperature, so degeneracy enters the transport
coefficients, and the thermal-diffusion and heat-flow coefficients must use
one consistent definition \cite{beznogov16,paxton18}. A coupled ion/electron
calculation on the saved gradients, including electron--electron collisions
and the leading ion recoil term, gives provisional outward hydrogen
velocities of $(1.101$--$1.303)\times10^{-9}$ cm s$^{-1}$ at enclosed mass
0.6972, with local drift times $H_P/|w_{\rm H}|$ of 12.78--15.13 Gyr there
and 3.694--4.244 Gyr near mass fraction 0.9738; the thermal force
contributes 23.63--24.48\% of the drift. The range spans two screening
choices and is not an uncertainty. Including electron--electron collisions
gives 0.5277--0.8780 of the elastic-only conductivity across the hot
nonconvective region. These are the coefficients Tracks S and F use above
about 2 MK; below that the tracks represent species exchange by convection
alone, and cold radiative layers still lack a supported transport law.
Heat carried by species motion is at most 0.09557\% of the local luminosity
on hot nonconvective faces; its omission below 2 MK changes surface
luminosity by at most 0.01320\% over 5 Myr.

\paragraph{Interior opacity sources.}
The hot interior opacity comes from one archive, and independent archives
disagree at the few-percent level. On 1940 states with $X=0$ or 0.05 and
$Z=0.01$ or 0.03, the Los Alamos OPLIB means distributed with MESA
\cite{farag24} differ from the TOPS means used here by a median of 2.364\%
and by up to 7.279\%; on 15 warm-envelope layers the Opacity Project
\cite{seaton05} differs from OPLIB by $-9.730$ to $+4.619\%$, and from TOPS
by 2.165\% at a sampled state. The archives use different element sets and
atomic physics, and no track has been run with the alternatives. Where the
sampled layers are convective or conductive the stellar response to opacity
changes of this size is unresolved (Table~\ref{tab:sensitivity}); in the
radiative envelope of the depleted star a $\pm10\%$ change moves the
effective temperature by 3.216 K, so a few-percent archive difference is
below the other uncertainties but is not zero.

\paragraph{The cold helium interior.}
The energy equation used throughout is
$L_\gamma=L_{\rm nuc,dep}+L_{\rm grav}-L_{\nu,\rm thermal}$ with nuclear
neutrinos removed from $L_{\rm nuc,dep}$; at $3.890\Tyr$ in Track P the
terms are 0.01077, $3.201\times10^{-5}$ and (nuclear neutrinos)
$3.362\times10^{-4}\Ls$, with plasma losses of $1.260\times10^{-7}\Ls$ and
a discrete residual below $8\times10^{-11}$. Beyond the computed tracks the
remnant needs a consistent liquid and solid helium free energy with quantum
and mixture effects \cite{baiko22}: specifying melting at $\Gamma=175$ does
not fix the latent heat or the liquid heat capacity. The cold-electron
thermal response, electron and ionic entropies and liquid/solid free
energies have been implemented and checked against thermodynamic identities
and against Ioffe EOSFI22 (largest density-response discrepancy
$1.05\times10^{-6}$ in ideal-ion pressure units), but they are not yet
combined into the production equation of state. Crystallization and its
latent heat, the depth at which convection couples to the isothermal core,
and pycnuclear rates in the actual helium-rich composition \cite{gasques05}
all remain to be evaluated.

\paragraph{Omitted reactions, neutrinos and mixing.}
Pep, hep, ppIII and the oxygen branches are absent in every track; adding a
closed CN cycle to a Track P state adds 2.381\% to the deposited pp energy
and lowers the surface luminosity by 0.1239\% after 5 Myr. Only plasma
neutrino losses are included. Semiconvective heat flux, thermohaline mixing
during the flash, rotation, magnetism, mass loss and accretion are absent.
The tracks describe an isolated, nonrotating star; accretion in particular
would change the surface composition on which the atmosphere depends
\cite{bauer19}.

\paragraph{External heating and encounters.}
Encounters, accretion and particle decay are treated in
Section~\ref{sec:lifetime} as conditional extensions of the isolated track
\cite{adams97}. For close encounters, energy transferred into stellar
oscillations is distinct from retained heat; gravitational radiation can
remove mode energy \cite{rathore05}, and the thermalized fraction must be
determined before it enters a heating history.
